\documentclass[12pt]{article}

\usepackage{graphicx}
\usepackage {pstcol,color,pst-node}
\usepackage{advi}
\usepackage{advi-annot}
\usepackage{advi-graphicx}
\usepackage{demo}
\usepackage{superpose}

\begin{document}

\newpage


\subsection* {\ActiveDVI~has two components}

\begin{itemize}
\item {\tt mldvi}: {\bf a dvi} previewer (by Alexandre Miquel)
\begin {itemize}
\item[+] dvips
  \textcolor{c1}{c}%
  \textcolor{c2}{o}%
  \textcolor{c3}{l}%
  \textcolor{c4}{o}%
  \textcolor{c5}{u}%
  \textcolor{c6}{r}
  extension (by Alexandre Miquel)
\item[+] psfile specials with cache (by Jun Furuse)
\item[+] GPIC specials (by Xavier Leroy)
\item[+] Postscript specials (by Didier R{\'{e}}my)
\item[+] hyperrefs (by Didier R{\'{e}}my)
\item[+] active anchors (by Didier R{\'{e}}my)
\item[+] background images (by Roberto Di Cosmo)
\item[+] and a few other gadgets...
\end {itemize}
\end{itemize}

\begin{itemize}
\item 
A set of macros to animate your presentation 
(by Jun~Furuse and Didier~R{\'{e}}my).
 
\item 
A simple way to launch arbitrary demonstration and/or animation
programs from inside your slides
\hbox {(by Jun Furuse and Pierre Weis)}.

\item 
A package to annotate text (using active anchors)

\item
Support for \verb+\includegraphics+ with native image formats

\end{itemize}

\newpage

%Didier's flashy macro
\def \flash #1{\let \do\leflash \do #1\relax}
\def \leflash #1{\ifx #1\relax \def \do{}\else \def \do {#1\adviwait[0.05]\leflash}\fi \do}

\def\goon{\hbox{{\Large $\Rightarrow$}~Type~space~to~go~on}}

\subsection* {EPS (and other!) graphics including using Camlimages}
\begin{quote}
\verb+\includegraphics[width=0.1\textwidth]{caml.eps}+\\
\verb+\includegraphics[width=0.2\textwidth]{caml.eps}+\\
\verb+\includegraphics[width=0.1\textwidth]{world}+\\

This last command looks for a file named {\tt world} with extensions out 
of the ones provided by the directive \verb+\DeclareGraphicsExtensions+
in the \verb|advi-images| package (in the present case, it finds \verb+world.jpg+, for which the
bounding box must be available in \verb+world.bb+)\\

\includegraphics[width=0.1\textwidth]{../tex/caml.eps}
\includegraphics[width=0.2\textwidth,height=0.09\textwidth]{../tex/caml.eps}
\framebox{\includegraphics[width=0.1\textwidth]{world}}
\end{quote}

\subsection* {Other eye candies for presentation}

\medskip\noindent
\docdef \adviwait\docopt{seconds}
\begin{quote}
Wait for \docid{seconds} or without argument forever, that is, until the
user requests to continue (or to change page).
\end{quote}

\medskip\noindent
\docdef \advirecord \doctt{[play]}\docarg{tag}\docarg{latex code}
\begin{quote}
The computes \docid{latex code} and record its DVI ouput under tag
\docid{tag}. The DVI output is not displayed during recording, unless the
option \doctt{play} is set.

Embedded recordings are prohibited. 
\end{quote}

\medskip\noindent
\docdef \adviplay\docopt{color}\docarg{tag}
\begin{quote}
Replay the DVI previously recorded as \docid{tag}.
The optional argument changes the color to \docid{color} during replay. 
\end{quote}


\newpage 

\subsection* {Eye candies (examples)}

\begin{itemize}
   \item mark the text with a tag name:\\
     $\tt \backslash advirecord\{this\}
      \{\advirecord[play]{this}{Warning~to~the~casual~user}\}$

   \item then change the current text color and replay the tagged text:\\
  $\tt{\backslash}play\{this\}$.
\end{itemize}


\noindent
For instance, to display in red the text associated to the tag
{\tt this} above, you can either pass a color option to the play command:
\verb+\adviplay[red]{this}}+\\
\advirecord[play]{that}{\goon}\adviwait%
\textcolor {c1}{\adviplay {this}\adviplay {that}}\adviwait
{\adviplay{this}\adviplay {that}}\adviwait

\noindent
or place the command in the appropriate color environment:
to display the same text in green, type\\
\verb+\textcolor{green}{\adviplay{this}}+\\
\advirecord[play]{goon}{\goon}\adviwait
\textcolor{green}{\adviplay{this}\adviplay{goon}\adviwait}

\noindent
Using pauses and color settings, you can get
\advirecord[play]{this}{various}
\textcolor{red}{\adviplay{this}\adviplay{goon}}\adviwait
\advirecord[play]{this}{color (and fonts)}
\textcolor{c2}{\adviplay{this}\adviplay{goon}}\adviwait
\advirecord[play]{this}{effects}
\textcolor{c3}{\adviplay{this}\adviplay{goon}}\adviwait
\advirecord[play]{this}{in}
\textcolor{c4}{\adviplay{this}\adviplay{goon}}\adviwait
\advirecord[play]{this}{your}
\textcolor{c5}{\adviplay{this}\adviplay{goon}}\adviwait
\advirecord[play]{this}{texts}
\textcolor{c6}{\adviplay{this}\adviplay{goon}}\adviwait
\advirecord[play]{this}{!}\\
\textcolor{red}{\adviplay{this}\adviplay{goon}}\adviwait
Including some strange flashring effect!
\adviplay{goon}
\textcolor{c6}{\adviplay{this}}\adviwait[\smallpause]%
\textcolor{c1}{\adviplay{this}}\adviwait[\smallpause]%
\textcolor{c2}{\adviplay{this}}\adviwait[\smallpause]%
\textcolor{c3}{\adviplay{this}}\adviwait[\smallpause]%
\textcolor{c4}{\adviplay{this}}\adviwait[\smallpause]%
\textcolor{c5}{\adviplay{this}}\adviwait[\smallpause]%
\textcolor{c6}{\adviplay{this}}\adviwait[\smallpause]%
\textcolor{c1}{\adviplay{this}}\adviwait[\smallpause]%
\textcolor{c2}{\adviplay{this}}\adviwait[\smallpause]%
\textcolor{c3}{\adviplay{this}}\adviwait[\smallpause]%
\textcolor{c4}{\adviplay{this}}\adviwait[\smallpause]%
\textcolor{c5}{\adviplay{this}}\adviwait[\smallpause]%
\textcolor{c6}{\adviplay{this}}\adviwait[\smallpause]%


\newpage

\subsection* {Embedded applications are easy to launch}

Any Unix command that obeys to the X geometry specifications may be
used to draw material in the {\ActiveDVI} window.
This is convenient to add a demo into your presentation.

\bigskip

\noindent
You launch a command using

\verb+\adviembed[width=w,height=h]{command -geometry !g -use !p}+

\noindent
where \verb"!g" specifies the required geometry, \verb"!p" the
{\ActiveDVI} window identifier, and \verb"w", \verb"h" the width \&
height of the rectangle where the application will draw its material.
Values of \verb"!p" and \verb"!g" are automatically computed by
{\ActiveDVI} and passed as arguments to the specified \verb"command".
 Examples:

\begin{itemize}
\item The \verb"display" command (from ImageMagick) can show any image into your
presentation using:\\
\verb|\adviembed[width=3cm,height=3.3cm]|\\
\verb|          {display -geometry !g! -window !p logo.gif}|\\
\adviembed[width=3cm, height=3.3260869cm]{display -geometry !g! -window !p advilogo.gif}

\item Similarly, \verb|animate| can display animations:\\
        \verb|\adviembed[width=2cm,height=2cm]|\\
        \verb|          {animate -geometry !g! -window !p anim.gif}|\\
        \adviembed[width=2cm,height=2cm]{animate -geometry !g! -window !p mmm.anim.gif}
\end{itemize}

% \newpage
% 
% \begin{center}
% {\large\bf I have four words for you!}\\[3mm]
% 
% \adviembed[width=10.62cm,height=7.92cm]{xanim +W!p +Sx!w +Sy!h +Zpe -Zr dancemonkeyboy.avi}
% \end{center}


\newpage

\subsection* {Embedded applications (cont.)}

\begin{itemize}
  \item Any application which has enbedding options like\\
   \verb+\adviembed[width=w,height=h]{prog -geometry !g -use !p}+\\
   For instance (using wish)\\
   \verb|\adviembed[persistent,width=5cm,height=1cm]|\\
   \verb|          {wish ./watch -geometry !g -use !p}|\\
   \adviembed[persistent,width=5cm,height=0.71cm]{wish ./watch -geometry !g -use !p}
\adviwait
  \item Possibly interactive ! (using CamlTk)\\
   \verb|\adviembed[persistent,width=6.5cm,height=9.5cm]|\\
   \verb|          {./taquin -geometry !g -use !p dojoji.gif}|\\
  \adviembed[persistent=taquin,width=7.5cm,height=11cm]{./taquin.sh -geometry !g -use !p dojoji.gif}

\end{itemize}

\newpage

\subsection* {Embedded applications (cont.)}

If we cannot specify the parent window of X applications,
they will be handled by the window manager.
Depending the WM, 
their placements may not obey your geometry specifications.
(Think about disabling ``remember window positions'' option of your WM.)

\begin{itemize}
 \item Xeyes using the macro command \verb|embed|:\\
   \verb|\adviembed[width=2cm,height=2cm]{xeyes -bw 0 -shape -geometry !g}|\\
   \adviembed[width=2cm,height=2cm]{xeyes -bw 0 -shape -geometry !g}

 \item An Xterm running an interesting application:\\
   \verb|\adviembed[width=10cm,height=5cm]{xterm -geometry 50x7+!x+!y -e ocaml}|\\
   \adviembed[width=10cm,height=4cm]{xterm -geometry 50x7+!x+!y -e ocaml}
\end{itemize}

The \verb|embed| macro can launch any executable system command during
the presentation. 
\begin{itemize}
 \item You can add music to your presentation, if you have a music
  player available:\\
  \verb|\adviembed[width=1mm,height=1mm]{mpg123 -q music.mp3}|
\end{itemize}

\advirecord{play}{\adviembed[name=mpg123,width=1mm,height=1mm]{mpg123 -q music.mp3}}
\advirecord {stop}{\advikillembed {mpg123}}
\advianchor[click]{play}{Play}
\advianchor[click]{stop}{Stop}

\newpage

\subsection* {Limitations of Postscript specials}

{\ActiveDVI} recognizes and renders Postscript Specials by calling and telling
\verb"ghostview" to display Postscript on the \verb"advi" window. 
Hence, images drawn using PS tricks will be displayed by \verb"advi". 

However, PS tricks make the assumptions that everything is displayed by
Postscript, which is still \verb"advi" displaying most of the material
itself. 

\subsubsection* {Most PS-tricks do work}
\begin {itemize}

\item[+]
Simple drawings work

\item[+] 
{\tt\string\SpecialCoor} works, {\em i.e.} commands of the form
{\tt \string \rput \{A\}\{bla bla\}} works where {\tt A} is a node

\item[+]
Connections between nodes {\tt \string \ncarc}, {\tt \string \ncarc},
also works.
\end {itemize}

\subsubsection* {Unfortunately, some PS-tricks do not work, yet}
\begin{itemize}
\item[-]
Labels over arrows \verb"\Aput", \verb"Bput", etc.
(they change the Postscript coordinates...)

\item[-]
{\tt pspicture}
(idem, but drawings that are not embedded in pspictures works)

\item[-]
Advi only recognizes binary overlays (display if 0, hide otherwise). 
This is sufficient for the proper documentclass
(see \hyperlink {overlays}{below}).

\end {itemize}



\subsection*{Known bug with PS tricks}

With some versions of \verb"gv" and/or \verb"X11", the PS tricks are not
correctly synchronized. Some drawings may hide text, or conversely.


% Examples: 

% \Rnode {A}{A} \qquad \Rnode {B}{B}
% \ncline{->}{A}{B} \Aput{L}

\newpage

\subsection* {Overlays}
\hypertarget {overlays}{}

Overlays of the \verb"seminar" mode uses PS tricks, but in a manner that
requires characters to be also printed as Postscript. The \verb"advi" package 
redefines \verb"\overlay" so that it draws only the first overlay (and
ignore others).

Hence, \verb"advi" cannot be used to preview slides using \verb"seminar"
overlays. However, a better package for overlays (build on top of seminar)
is \verb"prosper", which only uses the ability to show or hide some
material, which works well with \verb"advi".

Some reimplementation of \verb"\overlay" for \verb"advi" to work with 
seminar overlays might still be possible...  

\subsection* {Ignore mode}

The package \verb"advi" can be loaded with the \verb"batch" option:

\smallskip
{\verb"\usepackage[ignore]{advi}"}
\smallskip

The effect is to switch off some \verb"advi" gadgets and to render some others 
in Postscript, if possible. 


\newpage

\subsection* {Hyper links to other files}

\begin {itemize}

\item
(recursive) link to \hyperref{demo.dvi}{}{overlays}{\tt demo.dvi}
produced by

\qquad\verb"\hyperref{demo.dvi}{}{overlays}{\tt demo.dvi}"

A click opens file \verb"demo.dvi" with a new \verb"advi" and jumps to the
link \verb"overlays" in that file, as if the command

\qquad\verb"advi -html overlays demo.dvi"

had been typed in a shell.

(Currently, the file name should end with \verb".dvi" extension, but in the
future this should use Mine man Mailcap to find the type of the file and the
previewer for that file.)

% \item
% \hyperref{demo.tex}{}{foo}{\tt more demo.tex}

\item
Issue a warning if the extension is not recognized or if the 
file is non-existent:
\hyperref{nonexistent.dvi}{}{}{\tt nonexistent.dvi}
or
\hyperref{../tex/advi.sty}{}{}{\tt ../tex/advi.sty}



\end {itemize}


\newpage

\subsection* {Combination of replays and PS tricks}

Replay can be used to hide or show text or to change the color of text. 
However, postscript colors are hard-wired by {\TeX}, {\em i.e.} they are not
interpreted in the current color context as font colors. 

Replay can still be used to simulate replay of Postscript colors: instead, 
two identical figures can be superposed. For instance, 
$$
\psset {linestyle=none,fillstyle=solid}
\rlap
{\advirecord{tagb}{\ovalnode [fillcolor=green]{A}{\textbf {Postscript}}}}%
\advirecord[play]{taga}{\ovalnode [fillcolor=yellow]{B}{\textbf {Postscript}}}
$$
will  be rendered by
\begin{verbatim}
  \psset {linestyle=none,fillstyle=solid}
  \def \PS {\textbf {Postscript}}
  \rlap {\advirecord{tagb}{\ovalnode [fillcolor=green]{A}{\PS}}}%
  \advirecord[play]{taga}{\ovalnode [fillcolor=yellow]{B}{\PS}}}
\end{verbatim}
\begin {enumerate}

\item
Type space to change the text color. 
\adviwait \textcolor{red}{\adviplay{taga}}
\begin{verbatim}
    \adviwait \textcolor{red}{\adviplay{taga}}
\end{verbatim}

\item
Type space to change the background.
\adviwait\textcolor{red}{\adviplay{tagb}}
\begin{verbatim}
    \adviwait\textcolor{red}{\adviplay{tagb}}
\end{verbatim}

\end {enumerate}


\newpage

\subsection*{Active anchors}

Active anchors are much alike html anchors, except that instead of
jumping to a reference anchors (usually at another page), definitely, they 
allow to play a previously recorded piece of dvi when the mouse is over the
anchor region. The page is saved and reset to its normal shape when the mouse
leaves the anchor region. 

For instance, scratch those stamps with the mouse to see how much you win.
$$
\ovalnode [linestyle=none,fillstyle=solid,fillcolor=yellow]{A}
  {\advirecord {A}{\textbf {10000}}}
\advianchor{A}{\adviplay[yellow]{A}}
\hspace {4em}
\def \i{\textcolor{green}{.}}
\ovalnode [linestyle=none,fillstyle=solid,fillcolor=green]{B}
  {\advirecord {B}{\i\textbf{Lost}!\i}}
\advianchor{B}{\adviplay[green]{B}}
$$
For instance, the left game was typeset as
{\small
\begin{verbatim}
\ovalnode [linestyle=none,fillstyle=solid,fillcolor=yellow]{A}
  {\advirecord {A}{\textbf {10000}}}
\advianchor{A}{\adviplay[yellow]{A}}
\end{verbatim}}%

\medskip
\noindent
The effect of active anchors can also be 
\advianchor{C}{delocalized}.

\SpecialCoor
\advirecord {C}
  {\ncarc[linecolor=red,arcangleA=50,arcangleB=50,linewidth=3pt]
         {<->}{A}{B}%\Aput{Fire!}
}


\noindent
Another form of
\begin{superpose}
\advirecord[play]{a1}{\advianchor{a2}{simple}}\\
\advirecord{a2}{\adviplay[white]{a1}\colorbox{yellow}{ cool }}
\end{superpose}
simple annotations is the overriding of
\begin{superpose}
\advianchor{7}{words}\\
\advirecord{7}{\adviplay[white]{6}\colorbox{white}{{by text}}}\\
\end{superpose}
or by little 
\begin{superpose}
\advirecord[play]{1}{pieces \advianchor{2}{of} text}\\
\advirecord{2}{\colorbox{white}{\begin{tabular}{@{}l@{}}
inner tabular\\ should be between\\ parentheses\end{tabular}}}\\
\advirecord{3}{The \qquad End}
\end{superpose}

\newpage

\subsection*{\adviannot{Active Annotations}(2,1.5){Annotations \\are great!}}

\noindent
\adviannot[p=9em]{Active anchors}(1,1.5)
{\centering Active anchors are even greater!}
can also be used to annotate \adviannot[C,bg=green]{text}{text=t+e+x+t} as
here, \adviannot{for instance}[linecolor=red](-1.5,1){for example}. 

\medskip

\noindent
These annotations have been typeset respectively as:
\begin{small}
\begin{verbatim}
\adviannot{Active Annotations}(2,1.5){Annotations \\are great!}

\adviannot[p=9em]{Active anchors}(1,1.5)
   {\centering Active anchors are even greater!}

\adviannot[C,bg=green]{text}{text=+t+e+x+t}

\adviannot[C]{for instance}[linecolor=red](-1.5,1){for example?}
\end{verbatim}
\end{small}

\subsection*{Note}

The command 
\begin{quote}

\verb"\usepackage{advi}"

\verb"\usepackage{advi-annot}"

\end{quote}
must appear after other packages that uses PS tricks, since it redefines
some macros.  The package \verb"advi-annot" is necessary to get
the annotation macros.

\newpage

\subsection*{Superpose Package}

This very small package allows the superosing of boxes, making the
surrounding box so that it fits all included boxes. This is convenient
when displayin each included box at different pauses.


Movie...
\advirecord{0}{\colorbox{white}
{\begin{superpose}
\advirecord{1}{Begin}\\
\advirecord{2}{{\begin{tabular}{@{}l@{}}
inner tabular\\ should be between\\ parentheses\end{tabular}}}\\
\advirecord{3}{The End}
\end{superpose}}}
end
\adviplay{1}
\adviwait\adviplay{0}\adviplay{2}
\adviwait\adviplay{0}\adviplay{3}


\newpage

\subsection*{Pictures for \TeX}

{\ActiveDVI} can display pictures for \TeX, as generated by
the \verb"gpic" command into a list of \TeX~ \verb"\special" commands.

\bigskip

To obtain a figure, you first write commands for the language
{\em pic} into a file; from this pic file, you can generate a \TeX
file using a {\em pic compiler}, e.g. \verb"gpic"; then include this
compiled file into your presentation, and {\ActiveDVI} will nicely display it.

Let's consider the \verb"pic" source file \verb"trigo.pic" given
\hyperref{trigo_source.dvi}{}{}{here}; we obtain
\verb"trigo.tex" using the command 
\begin{verbatim}
gpic -t trigo.pic > trigo.tex
\end{verbatim}

Then we include \verb"trigo.tex" with the \TeX command
\verb|\expandafter\ifx\csname graph\endcsname\relax \csname newbox\endcsname\graph\fi
\expandafter\ifx\csname graphtemp\endcsname\relax \csname newdimen\endcsname\graphtemp\fi
\setbox\graph=\vtop{\vskip 0pt\hbox{%
    \special{pn 8}%
    \special{ar 1050 1350 900 900 0 6.28319}%
    \graphtemp=.5ex\advance\graphtemp by 1.440in
    \rlap{\kern 0.960in\lower\graphtemp\hbox to 0pt{\hss O\hss}}%
    \graphtemp=.5ex\advance\graphtemp by 2.550in
    \rlap{\kern 1.050in\lower\graphtemp\hbox to 0pt{\hss The trigonometric circle\hss}}%
    \graphtemp=\baselineskip\multiply\graphtemp by 1\divide\graphtemp by 2
    \advance\graphtemp by .5ex\advance\graphtemp by 1.350in
    \rlap{\kern 1.950in\lower\graphtemp\hbox to 0pt{~1\hss}}%
    \graphtemp=\baselineskip\multiply\graphtemp by -1\divide\graphtemp by 2
    \advance\graphtemp by .5ex\advance\graphtemp by 0.450in
    \rlap{\kern 1.050in\lower\graphtemp\hbox to 0pt{\hss 1~}}%
    \special{pa 0 1350}%
    \special{pa 2400 1350}%
    \special{fp}%
    \special{sh 1.000}%
    \special{pa 2300 1325}%
    \special{pa 2400 1350}%
    \special{pa 2300 1375}%
    \special{pa 2300 1325}%
    \special{fp}%
    \graphtemp=\baselineskip\multiply\graphtemp by 1\divide\graphtemp by 2
    \advance\graphtemp by .5ex\advance\graphtemp by 1.350in
    \rlap{\kern 2.400in\lower\graphtemp\hbox to 0pt{\hss $x$}}%
    \special{pa 1050 2400}%
    \special{pa 1050 0}%
    \special{fp}%
    \special{sh 1.000}%
    \special{pa 1025 100}%
    \special{pa 1050 0}%
    \special{pa 1075 100}%
    \special{pa 1025 100}%
    \special{fp}%
    \graphtemp=.5ex\advance\graphtemp by 0.000in
    \rlap{\kern 1.050in\lower\graphtemp\hbox to 0pt{\hss $y$~}}%
    \special{ar 1282 1319 250 250 -0.886486 0.125474}%
    \special{sh 1.000}%
    \special{pa 1502 1208}%
    \special{pa 1440 1125}%
    \special{pa 1533 1169}%
    \special{pa 1502 1208}%
    \special{fp}%
    \graphtemp=\baselineskip\multiply\graphtemp by -1\divide\graphtemp by 2
    \advance\graphtemp by .5ex\advance\graphtemp by 1.319in
    \rlap{\kern 1.532in\lower\graphtemp\hbox to 0pt{~$\theta$\hss}}%
    \special{pa 1050 1350}%
    \special{pa 1829 900}%
    \special{fp}%
    \graphtemp=.5ex\advance\graphtemp by 0.800in
    \rlap{\kern 1.929in\lower\graphtemp\hbox to 0pt{\hss $M$\hss}}%
    \special{pa 1829 900}%
    \special{pa 1050 900}%
    \special{fp}%
    \special{sh 1.000}%
    \special{pa 1150 925}%
    \special{pa 1050 900}%
    \special{pa 1150 875}%
    \special{pa 1150 925}%
    \special{fp}%
    \graphtemp=.5ex\advance\graphtemp by 1.125in
    \rlap{\kern 0.810in\lower\graphtemp\hbox to 0pt{\hss $\sin\left(\theta\right)\left\{\begin{array}{c} \\ \\ \end{array}\right.$\hss}}%
    \special{pa 1829 900}%
    \special{pa 1829 1350}%
    \special{fp}%
    \special{sh 1.000}%
    \special{pa 1854 1250}%
    \special{pa 1829 1350}%
    \special{pa 1804 1250}%
    \special{pa 1854 1250}%
    \special{fp}%
    \graphtemp=.5ex\advance\graphtemp by 1.350in
    \rlap{\kern 1.829in\lower\graphtemp\hbox to 0pt{\hss $\underbrace{\hspace*{1.95cm}}_{\displaystyle\cos~\left(\theta\right)~}$}}%
    \hbox{\vrule depth2.550in width0pt height 0pt}%
    \kern 2.400in
  }%
}%
\showgraph
| and get:

\def\showgraph{%
  \par\medskip\centerline{\raise 1em\box\graph}\bigskip\noindent\ignorespaces}

\expandafter\ifx\csname graph\endcsname\relax \csname newbox\endcsname\graph\fi
\expandafter\ifx\csname graphtemp\endcsname\relax \csname newdimen\endcsname\graphtemp\fi
\setbox\graph=\vtop{\vskip 0pt\hbox{%
    \special{pn 8}%
    \special{ar 1050 1350 900 900 0 6.28319}%
    \graphtemp=.5ex\advance\graphtemp by 1.440in
    \rlap{\kern 0.960in\lower\graphtemp\hbox to 0pt{\hss O\hss}}%
    \graphtemp=.5ex\advance\graphtemp by 2.550in
    \rlap{\kern 1.050in\lower\graphtemp\hbox to 0pt{\hss The trigonometric circle\hss}}%
    \graphtemp=\baselineskip\multiply\graphtemp by 1\divide\graphtemp by 2
    \advance\graphtemp by .5ex\advance\graphtemp by 1.350in
    \rlap{\kern 1.950in\lower\graphtemp\hbox to 0pt{~1\hss}}%
    \graphtemp=\baselineskip\multiply\graphtemp by -1\divide\graphtemp by 2
    \advance\graphtemp by .5ex\advance\graphtemp by 0.450in
    \rlap{\kern 1.050in\lower\graphtemp\hbox to 0pt{\hss 1~}}%
    \special{pa 0 1350}%
    \special{pa 2400 1350}%
    \special{fp}%
    \special{sh 1.000}%
    \special{pa 2300 1325}%
    \special{pa 2400 1350}%
    \special{pa 2300 1375}%
    \special{pa 2300 1325}%
    \special{fp}%
    \graphtemp=\baselineskip\multiply\graphtemp by 1\divide\graphtemp by 2
    \advance\graphtemp by .5ex\advance\graphtemp by 1.350in
    \rlap{\kern 2.400in\lower\graphtemp\hbox to 0pt{\hss $x$}}%
    \special{pa 1050 2400}%
    \special{pa 1050 0}%
    \special{fp}%
    \special{sh 1.000}%
    \special{pa 1025 100}%
    \special{pa 1050 0}%
    \special{pa 1075 100}%
    \special{pa 1025 100}%
    \special{fp}%
    \graphtemp=.5ex\advance\graphtemp by 0.000in
    \rlap{\kern 1.050in\lower\graphtemp\hbox to 0pt{\hss $y$~}}%
    \special{ar 1282 1319 250 250 -0.886486 0.125474}%
    \special{sh 1.000}%
    \special{pa 1502 1208}%
    \special{pa 1440 1125}%
    \special{pa 1533 1169}%
    \special{pa 1502 1208}%
    \special{fp}%
    \graphtemp=\baselineskip\multiply\graphtemp by -1\divide\graphtemp by 2
    \advance\graphtemp by .5ex\advance\graphtemp by 1.319in
    \rlap{\kern 1.532in\lower\graphtemp\hbox to 0pt{~$\theta$\hss}}%
    \special{pa 1050 1350}%
    \special{pa 1829 900}%
    \special{fp}%
    \graphtemp=.5ex\advance\graphtemp by 0.800in
    \rlap{\kern 1.929in\lower\graphtemp\hbox to 0pt{\hss $M$\hss}}%
    \special{pa 1829 900}%
    \special{pa 1050 900}%
    \special{fp}%
    \special{sh 1.000}%
    \special{pa 1150 925}%
    \special{pa 1050 900}%
    \special{pa 1150 875}%
    \special{pa 1150 925}%
    \special{fp}%
    \graphtemp=.5ex\advance\graphtemp by 1.125in
    \rlap{\kern 0.810in\lower\graphtemp\hbox to 0pt{\hss $\sin\left(\theta\right)\left\{\begin{array}{c} \\ \\ \end{array}\right.$\hss}}%
    \special{pa 1829 900}%
    \special{pa 1829 1350}%
    \special{fp}%
    \special{sh 1.000}%
    \special{pa 1854 1250}%
    \special{pa 1829 1350}%
    \special{pa 1804 1250}%
    \special{pa 1854 1250}%
    \special{fp}%
    \graphtemp=.5ex\advance\graphtemp by 1.350in
    \rlap{\kern 1.829in\lower\graphtemp\hbox to 0pt{\hss $\underbrace{\hspace*{1.95cm}}_{\displaystyle\cos~\left(\theta\right)~}$}}%
    \hbox{\vrule depth2.550in width0pt height 0pt}%
    \kern 2.400in
  }%
}%
\showgraph


See the \verb"gpic" command manual for more information about
\verb"gpic" and the programming language \verb"pic".

\lastpage

~\vfill
\begin{center}
\includegraphics[width=0.5\textwidth]{../tex/advilogo.eps}\\
{\Large \bf That's all, folks!!}
\end{center}
\vfill

\end{document}
